% GENERATED FILE. Edit canonical lessons, then run npm run generate:books.
% canonical-source-hash: fnv1a64:d43044a548e30ac6
% canonical-lessons: GE-C230-gegenteil, GE-C230-geschichte, GE-C230-gespraech, GE-C230-gewicht, GE-C230-gewohnheit

\chapter{The opposite, A story, A conversation, Weight, A habit}
\label{ch:ge-the-opposite-a-story-a-conversation-weight-a-habit}
\hlchaptermodality{230}

\begin{chapteropening}
\textbf{By the end of this chapter:} \emph{I can say the opposite, a story, a conversation, weight and a habit.}

Complete the last lesson of chapter 230: I can say the opposite, a story, a conversation, weight and a habit.
\end{chapteropening}

\section[das Gegenteil]{das Gegenteil — the opposite}

\label{lesson:GE-C230-gegenteil}



\begin{quote}
Before the new one: say the German for a feeling, then the German for an area.
\end{quote}

\subsection*{You'll want to know: das Gegenteil}
\textbf{das Gegenteil} — \textquotedblleft{}the opposite\textquotedblright{}.

\begin{cousinweb}
\textbf{Gegenteil} is the word for \textquotedblleft{}against\textquotedblright{} + the word for \textquotedblleft{}part\textquotedblright{}, Old High German \emph{teil} $\leftarrow$ Proto-Germanic \emph{*dailą}, English \emph{deal} (EtymDB 2.1): the opposite part.
\end{cousinweb}

\begin{grammarlens}[title={What you've built}]
One more word for this chapter.
\end{grammarlens}

\subsection*{Your turn}
\begin{itemize}
\raggedright
  \item \emph{Say it:} \emph{das Gegenteil}
  \item \emph{Say it:} \emph{das Gegenteil}, once more
  \item \emph{Say it:} say \emph{das Gegenteil}
  \item \emph{Recall:} say the German for a feeling, then the German for an area, then say \emph{das Gegenteil} again
\end{itemize}

\begin{tcolorbox}[breakable,colback=teal!4,colframe=teal!35!black,title={Before you move on}]
What does das Gegenteil mean? (\textquotedblleft{}The opposite\textquotedblright{}.) Say it once more.
\end{tcolorbox}
\section[die Geschichte]{die Geschichte — a story, history}

\label{lesson:GE-C230-geschichte}



\begin{quote}
Before the new one: say the German for an area, then the German for the opposite.
\end{quote}

\subsection*{You'll want to know: die Geschichte}
\textbf{die Geschichte} — \textquotedblleft{}a story, history\textquotedblright{}.

\begin{cousinweb}
\textbf{Geschichte} $\leftarrow$ Middle High German \emph{geschicht} $\leftarrow$ Old High German \emph{gisciht}, \textquotedblleft{}event\textquotedblright{} (EtymDB 2.1) — inherited, from the verb \textquotedblleft{}to happen\textquotedblright{}. A story is what happened, told.
\end{cousinweb}

\begin{grammarlens}[title={What you've built}]
One more word for this chapter.
\end{grammarlens}

\subsection*{Your turn}
\begin{itemize}
\raggedright
  \item \emph{Say it:} \emph{die Geschichte}
  \item \emph{Say it:} \emph{die Geschichte}, once more
  \item \emph{Say it:} say \emph{die Geschichte}
  \item \emph{Recall:} say the German for an area, then the German for the opposite, then say \emph{die Geschichte} again
\end{itemize}

\begin{tcolorbox}[breakable,colback=teal!4,colframe=teal!35!black,title={Before you move on}]
What does die Geschichte mean? (\textquotedblleft{}A story, history\textquotedblright{}.) Say it once more.
\end{tcolorbox}
\section[das Gespräch]{das Gespräch — a conversation}

\label{lesson:GE-C230-gespraech}



\begin{quote}
Before the new one: say the German for the opposite, then the German for a story.
\end{quote}

\subsection*{You'll want to know: das Gespräch}
\textbf{das Gespräch} — \textquotedblleft{}a conversation\textquotedblright{}.

\begin{cousinweb}
\textbf{Gespräch} $\leftarrow$ Middle High German \emph{gespræche} $\leftarrow$ Old High German \emph{gisprāhhi} (EtymDB 2.1): \emph{ge-}, \textquotedblleft{}together\textquotedblright{}, + the root of \emph{Sprache}. A conversation is speech together.
\end{cousinweb}

\begin{grammarlens}[title={What you've built}]
One more word for this chapter.
\end{grammarlens}

\subsection*{Your turn}
\begin{itemize}
\raggedright
  \item \emph{Say it:} \emph{das Gespräch}
  \item \emph{Say it:} \emph{das Gespräch}, once more
  \item \emph{Say it:} say \emph{das Gespräch}
  \item \emph{Recall:} say the German for the opposite, then the German for a story, then say \emph{das Gespräch} again
\end{itemize}

\begin{tcolorbox}[breakable,colback=teal!4,colframe=teal!35!black,title={Before you move on}]
What does das Gespräch mean? (\textquotedblleft{}A conversation\textquotedblright{}.) Say it once more.
\end{tcolorbox}
\section[das Gewicht]{das Gewicht — weight}

\label{lesson:GE-C230-gewicht}



\begin{quote}
Before the new one: say the German for a story, then the German for a conversation.
\end{quote}

\subsection*{You'll want to know: das Gewicht}
\textbf{das Gewicht} — \textquotedblleft{}weight\textquotedblright{}.

\begin{cousinweb}
\textbf{Gewicht} is \emph{ge-} + the root of \emph{wiegen}, \textquotedblleft{}to weigh\textquotedblright{} — EtymDB 2.1 lists it beside English \emph{weight} and Dutch \emph{gewicht}.
\end{cousinweb}

\begin{grammarlens}[title={What you've built}]
One more word for this chapter.
\end{grammarlens}

\subsection*{Your turn}
\begin{itemize}
\raggedright
  \item \emph{Say it:} \emph{das Gewicht}
  \item \emph{Say it:} \emph{das Gewicht}, once more
  \item \emph{Say it:} say \emph{das Gewicht}
  \item \emph{Recall:} say the German for a story, then the German for a conversation, then say \emph{das Gewicht} again
\end{itemize}

\begin{tcolorbox}[breakable,colback=teal!4,colframe=teal!35!black,title={Before you move on}]
What does das Gewicht mean? (\textquotedblleft{}Weight\textquotedblright{}.) Say it once more.
\end{tcolorbox}
\section[die Gewohnheit]{die Gewohnheit — a habit}

\label{lesson:GE-C230-gewohnheit}



\begin{quote}
Before the new one: say the German for a conversation, then the German for weight.
\end{quote}

\subsection*{You'll want to know: die Gewohnheit}
\textbf{die Gewohnheit} — \textquotedblleft{}a habit\textquotedblright{}.

\begin{cousinweb}
\textbf{Gewohnheit} is \emph{gewohnt}, \textquotedblleft{}accustomed\textquotedblright{}, + \emph{-heit}: what you are accustomed to. \emph{gewohnt} is built on \emph{wohnen}, first \textquotedblleft{}to be accustomed, to stay\textquotedblright{}.
\end{cousinweb}

\begin{grammarlens}[title={What you've built}]
One more word for this chapter.
\end{grammarlens}

\subsection*{Your turn}
\begin{itemize}
\raggedright
  \item \emph{Say it:} \emph{die Gewohnheit}
  \item \emph{Say it:} \emph{die Gewohnheit}, once more
  \item \emph{Say it:} say \emph{die Gewohnheit}
  \item \emph{Recall:} say the German for a conversation, then the German for weight, then say \emph{die Gewohnheit} again
\end{itemize}

\begin{tcolorbox}[breakable,colback=teal!4,colframe=teal!35!black,title={Before you move on}]
What does die Gewohnheit mean? (\textquotedblleft{}A habit\textquotedblright{}.) Say it once more.
\end{tcolorbox}
